% ID: MAT-ALG-INT-012
% Titolo: Espressione con potenze dispari e basi negative
% Disciplina: Matematica
% Area: algebra
% Argomento: numeri_interi
% Classe: Prima secondaria di secondo grado
% Difficolta: 3
% Tipo: problema_numerico
% Tag: matematica, algebra, numeri_interi
% Fonte: appunti_matematica_anonimizzati
% Provenienza: materiale fornito dal docente; testo anonimizzato; soluzione redatta per Exergo

\begin{esercizio}
Calcola $(-5^5)^5:\left[(-125)^2\cdot(-25)^3\right]^2+\left[(-5)^5\right]^5:(-125)^8$.
\end{esercizio}

\begin{soluzione}
Il primo numeratore è $-5^{25}$; il suo denominatore è $[5^6(-5^6)]^2=5^{24}$, quindi il primo quoziente è $-5$. Il secondo quoziente è $(-5)^{25}/5^{24}=-5$. Somma: $-10$.
\end{soluzione}
